\subsection{Population priors and the relation to TwinEB}
\label{sec:twineb}
The closest comparison for flexible empirical-Bayes prior learning is
TwinEB, introduced by Salehi, Nazaret, Shah and Blei
\cite{twineb}. TwinEB learns mixture priors on the entire row and column
latent vectors and fits these priors jointly with a variational posterior.
The paper develops Gaussian and Poisson matrix factorization.
Its mixture-of-vectors construction must be distinguished from a product
of independent scalar mixture priors: a shared component label can induce
dependence among coordinates of a latent vector.

For the Gaussian formulation, write a diagonal-component row mixture as
\begin{equation}\label{eq:twinebvector}
p_{\rm pop}(z)=\sum_{h=1}^H\omega_h
       \prod_{k=1}^K\Normal{z_k}{m_{hk}}{s_{hk}^2}.
\end{equation}
The scales in the paper are parameterized per component and coordinate
(Appendix B). Its mean-field posterior approximation does not imply
independence under the prior.
For example, with \(\bar m=\sum_h\omega_hm_h\),
\[
\Cov(z)=\sum_h\omega_h\operatorname{diag}(s_h^2)
       +\sum_h\omega_h(m_h-\bar m)(m_h-\bar m)^\top.
\]
The second term may have nonzero off-diagonal entries. Thus dependence
between latent programs is not, by itself, a capability absent from TwinEB.

\paragraph{An exact relationship at the level of the loading prior.}
The vector-mixture form implies the autoregressive conditionals
\begin{align}
p_{\rm pop}(z_k\mid z_{<k})
&=\sum_h\widetilde\omega_{hk}(z_{<k})
              \Normal{z_k}{m_{hk}}{s_{hk}^2},\label{eq:twinebconditional}\\
\widetilde\omega_{hk}(z_{<k})
&=\frac{\omega_h\prod_{j<k}\Normal{z_j}{m_{hj}}{s_{hj}^2}}
        {\sum_r\omega_r\prod_{j<k}\Normal{z_j}{m_{rj}}{s_{rj}^2}}.
\label{eq:twinebgate}
\end{align}
To prove the result, let
\(C_k(z_{\leq k})=\sum_h\omega_h
\prod_{j\leq k}\Normal{z_j}{m_{hj}}{s_{hj}^2}\), with \(C_0=1\).
The conditional in \eqref{eq:twinebconditional} equals
\(C_k/C_{k-1}\). Multiplying over \(k\) telescopes to \(C_K\), which
is \eqref{eq:twinebvector}.

Accordingly, a conditional Gaussian-mixture framework that permits the
specific gates \eqref{eq:twinebgate} contains this row-prior family exactly.
An unrestricted conditional head can also vary the component means and
variances with preceding coordinates, and incorporate observed covariates.
However, an arbitrary finite ReLU network does not necessarily implement
the log-quadratic Bayes gates exactly. A one-slab CGB is also a much
smaller family than an unrestricted vector mixture.
The result is a distributional inclusion under explicit assumptions,
not a blanket inclusion claim for every registry configuration.

Any joint density admits an autoregressive factorization by the chain
rule. Moreover, increasing the number of Gaussian mixture components
already gives highly flexible continuous population models.
The proposed ordering therefore does not establish that TwinEB cannot
represent nonlinear dependence. The empirical question is whether the
chosen conditional parameterization represents and learns scientifically
relevant dependencies more efficiently at feasible sample sizes and
computational budgets. Exact atomic probabilities are an additional
modeling choice; narrow continuous Gaussians approximate concentration
near zero without assigning positive probability to zero itself.

\paragraph{Different structural assumptions.}
The present model explicitly separates modality-specific latent programs
and specifies
\[
p(A_i,R_i\mid X_i)=p(A_i\mid X_i)p(R_i\mid A_i,X_i)
\]
through within- and between-modality conditional distributions.
This allows external covariates to modify activation probabilities,
amplitudes or scales at particular nodes. All such conditional factors
belong to one normalized joint model; upstream updates include the
downstream evidence described in \eqref{eq:familyD}.

TwinEB, as formulated in the cited paper, uses population priors for a
single matrix. A multi-view extension could concatenate loading vectors
and specify a suitable block observation model. That is a legitimate
comparator construction, but it should be identified as an extension
and implemented with the same missing-data likelihood and evaluation
protocol rather than attributed to the published model.

The feature-factor assumptions differ as well. TwinEB's column prior is
a mixture on an entire feature's latent vector; it can relate factor
coordinates within that vector. The present HMM instead relates adjacent
feature positions within each factor and is independent across factors.
It is not a strict superset of the TwinEB column-prior model. Similarly,
the current Gaussian observation implementation does not subsume the
paper's Poisson formulation.

\paragraph{Inference and the proposed contribution.}
Both approaches use learned distributions to regularize matrix
factorization. TwinEB optimizes prior and mean-field variational
parameters through an ELBO. The reference solver here uses analytic
normal-means proposals with a child-likelihood Metropolis correction,
sequence sampling for HMM factors, and optional partial Monte Carlo
empirical-Bayes learning. Section~\ref{sec:familyvi} supplies an alternative
augmented variational formulation; it has not been established to be
faster or more accurate than TwinEB's inference.

A defensible methodological contribution is therefore a modular
conditional-prior factorization framework with explicit sparse or
positive distributions, fixed side information, ordered feature priors,
and joint uncertainty propagation across partially observed modalities.
The distributional decomposition and the Metropolis identity are
established mathematics. Novelty would rest on a useful model
construction, effective inference and evidence that these structural
choices improve a defined scientific task.

\paragraph{Comparisons needed to support this positioning.}
Include TwinEB alongside independent cEBMF and the multi-omics
comparators below. On a common Gaussian single-matrix benchmark,
compare vector population mixtures and autoregressive mixtures at
matched ranks and documented parameter and tuning budgets. For
multi-view data, report any required extension of the competitor.
Separately ablate the hierarchy, nonlinear means and scales, exact
spikes, exponential slabs, fixed covariates and HMM smoothing.
This distinguishes prior expressiveness, regularization, modality
coupling and the inference algorithm.

Evaluation should include held-out observations, missing-modality
prediction, predictive calibration, and computation at matched
accuracy. For samplers, report effective sample sizes and convergence
diagnostics rather than acceptance rates alone. For VI, assess the
effect of its dependence assumptions against small reference problems.
No result in this derivation establishes empirical superiority over
TwinEB or the other methods.
